\textbf{Established partial result.}
Put \(x_n(c)=n^{4/3}d_n(c)\), \(\alpha_c=C_0(c)\kappa_c\), and \(\beta_c=C_0(c)C_A\).  The universal second-order rate theorem gives
\[
 0<\alpha_c\leq\liminf_{n\to\infty}x_n(c)
 \leq\limsup_{n\to\infty}x_n(c)\leq\beta_c<\infty.
\]
Consequently, the subsequential-limit set
\[
 \mathcal L(c)=\bigcap_{N\geq1}
 \overline{\{x_n(c):n\geq N\}}
\]
is a nonempty compact subset of \([\alpha_c,\beta_c]\), and all its elements are finite and strictly positive.

There is an exact criterion for a slowly varying correction.  A regularly varying normalizer satisfying \(a_nd_n(c)\to C\in(0,\infty)\) must have the form \(a_n=n^{4/3}\ell_n\), with \(\ell_n\) slowly varying and eventually bounded above and away from zero.  Such \(\ell_n\) and \(C\) exist if and only if \((x_n(c))\) is itself slowly varying.  Indeed, if \(\ell_nx_n(c)\to C\), then for every \(s>0\),
\[
 \frac{x_{\lfloor sn\rfloor}(c)}{x_n(c)}
 =\frac{\ell_n}{\ell_{\lfloor sn\rfloor}}
 \frac{\ell_{\lfloor sn\rfloor}x_{\lfloor sn\rfloor}(c)}
 {\ell_nx_n(c)}\longrightarrow1.
\]
Conversely, eventual positivity and slow variation of \(x_n(c)\) permit \(\ell_n=C/x_n(c)\), which is slowly varying and makes the product equal to \(C\) eventually.  Thus a nonconvergent normalized risk can be repaired by a slowly varying factor exactly when its nonconvergence remains within the slowly varying class.

The certificate conclusion now holds for every rational nonzero zero-sum contrast, not only for \(c^\dagger\).  For such \(c\), choose \(M_n\) with \(n^{4/3}/M_n^2\to0\), for example \(M_n=n\), and define
\[
 \begin{aligned}
 J_n^{\Lambda}(c)&=n^{4/3}\{C_0(c)n^{-1}-\Lambda_{n,M_n}(c)\},\\
 L_n^{\mathrm{cert}}(c)&=n^{4/3}\{C_0(c)n^{-1}-U_{n,M_n}(c)\},\\
 U_n^{\mathrm{cert}}(c)&=n^{4/3}\{C_0(c)n^{-1}-B_{n,c}(\nu_{n,M_n,c})\}.
 \end{aligned}
\]
The rational-contrast certificate theorem gives
\[
 J_n^{\Lambda}(c)\leq L_n^{\mathrm{cert}}(c)
 \leq x_n(c)\leq U_n^{\mathrm{cert}}(c),
 \qquad
 0\leq U_n^{\mathrm{cert}}(c)-J_n^{\Lambda}(c)
 \leq\frac{C_0(c)n^{4/3}}{4M_n^2}\longrightarrow0.
\]
Hence all four sequences differ by \(o(1)\), have exactly the same subsequential-limit set, and converge simultaneously to the same value if any one converges.  Because they are eventually bounded away from zero, slow variation of any one is equivalent to slow variation of all.  The contrast \(c^\dagger\) is the computed showcase: its existing exact rational calculations through \(n=5\) instantiate these general certificates, but they do not prove convergence, select a subsequential limit, or determine a limiting operator.